
\documentclass{article}
\usepackage{colortbl}
\usepackage{makecell}
\usepackage{multirow}
\usepackage{supertabular}

\begin{document}

\newcounter{utterance}

\centering \large Interaction Transcript for game `dond', experiment `coop', episode 1 with mistral{-}large{-}3{-}azure.
\vspace{24pt}

{ \footnotesize  \setcounter{utterance}{1}
\setlength{\tabcolsep}{0pt}
\begin{supertabular}{c@{$\;$}|p{.15\linewidth}@{}p{.15\linewidth}p{.15\linewidth}p{.15\linewidth}p{.15\linewidth}p{.15\linewidth}}
    \# & $\;$A & \multicolumn{4}{c}{Game Master} & $\;\:$B\\
    \hline

    \theutterance \stepcounter{utterance}  
    & & \multicolumn{4}{p{0.6\linewidth}}{
        \cellcolor[rgb]{0.9,0.9,0.9}{
            \makecell[{{p{\linewidth}}}]{
                \texttt{\tiny{[P1$\langle$GM]}}
                \texttt{Pelaat yhteistoiminnallista neuvottelupeliä, jossa sinun on sovittava toisen pelaajan kanssa siitä, kuinka jakaa joukko kohteita teidän välillä.} \\
\\ 
\texttt{Säännöt:} \\
\texttt{(a) Sinä ja toinen pelaaja saatte joukon kohteita. Kummallekin annetaan myös salainen arvofunktio, joka kuvaa, kuinka paljon arvostit kunkin kohteen tyyppiä.} \\
\texttt{(b) Vaihdat viestejä toisen pelaajan kanssa sopimaan siitä, kuka saa mitkä kohteet. Jokainen voi lähettää enintään 5 viestiä tai lopettaa pelin aikaisin tekemällä salaisen ehdotuksen milloin tahansa.} \\
\texttt{(c) Sinua pyydetään lähettämään salainen ehdotus, joka osoittaa kohteet, jotka haluat, muotoiltuna hakasulkeiden sisällä seuraavasti: "{[}Ehdotus: <numero> <kohteen nimi>, <numero> <kohteen nimi>, <...>{]}"} \\
\texttt{(d) Jos ehdotuksesi ovat täydentäviä, eli molempien ehdotusten täyttämiseksi on tarpeeksi kohteita, jokainen pelaaja saa pistemäärän saamiensa kohteiden arvojen summan perusteella. Muussa tapauksessa molemmat saavat nolla pistettä.} \\
\texttt{(e) Tavoitteesi on maksimoida pistemääräsi ja toisen pelaajan pistemäärän summa.} \\
\\ 
\texttt{Aloitetaan.} \\
\\ 
\texttt{Saatavilla olevien kohteiden joukko on:} \\
\\ 
\texttt{2 kynät, 2 sitruunat, 2 potit} \\
\\ 
\texttt{Salainen arvofunktiosi on:} \\
\\ 
\texttt{kynä: 2, sitruuna: 0, potti: 3} \\
\\ 
\texttt{TÄRKEÄÄ: Viestit, paitsi salaiset ehdotukset, lähetetään suoraan toiselle pelaajalle, joten älä sisällytä mitään vastausta sääntöihin tai tekstiä, joka ilmoittaa viestistäsi. Salaisesti ehdotuksen tekemiseksi käytä ilmoitettua muotoa. Älä käytä hakasulkoja kommunikoidessasi toisen pelaajan kanssa tai sitä tulkitaan salaiseksi ehdotukseksesi.} \\
            }
        }
    }
    & \\ \\

    \theutterance \stepcounter{utterance}  
    & \multicolumn{4}{p{0.6\linewidth}}{
        \cellcolor[rgb]{0.9,0.9,0.9}{
            \makecell[{{p{\linewidth}}}]{
                \texttt{\tiny{[P1$\rangle$GM]}}
                \texttt{Minulla on kynät ja potit arvokkaimpia. Ehdotan, että otan yhden potin ja yhden kynän, jos sinulle sopii. Miten sinä näet kohteiden arvot?} \\
            }
        }
    }
    & & \\ \\

    \theutterance \stepcounter{utterance}  
    & & & \multicolumn{2}{p{0.3\linewidth}}{
        \cellcolor[rgb]{0.9,0.9,0.9}{
            \makecell[{{p{\linewidth}}}]{
                \texttt{\tiny{[GM$|$GM]}}
                \texttt{continue} \\
            }
        }
    }
    & & \\ \\

    \theutterance \stepcounter{utterance}  
    & & & \multicolumn{2}{p{0.3\linewidth}}{
        \cellcolor[rgb]{0.9,0.9,0.9}{
            \makecell[{{p{\linewidth}}}]{
                \texttt{\tiny{[GM$|$GM]}}
                \texttt{Minulla on kynät ja potit arvokkaimpia. Ehdotan, että otan yhden potin ja yhden kynän, jos sinulle sopii. Miten sinä näet kohteiden arvot?} \\
            }
        }
    }
    & & \\ \\

    \theutterance \stepcounter{utterance}  
    & & \multicolumn{4}{p{0.6\linewidth}}{
        \cellcolor[rgb]{0.9,0.9,0.9}{
            \makecell[{{p{\linewidth}}}]{
                \texttt{\tiny{[GM$\rangle$P2]}}
                \texttt{Pelaat yhteistoiminnallista neuvottelupeliä, jossa sinun on sovittava toisen pelaajan kanssa siitä, kuinka jakaa joukko kohteita teidän välillä.} \\
\\ 
\texttt{Säännöt:} \\
\texttt{(a) Sinä ja toinen pelaaja saatte joukon kohteita. Kummallekin annetaan myös salainen arvofunktio, joka kuvaa, kuinka paljon arvostit kunkin kohteen tyyppiä.} \\
\texttt{(b) Vaihdat viestejä toisen pelaajan kanssa sopimaan siitä, kuka saa mitkä kohteet. Jokainen voi lähettää enintään 5 viestiä tai lopettaa pelin aikaisin tekemällä salaisen ehdotuksen milloin tahansa.} \\
\texttt{(c) Sinua pyydetään lähettämään salainen ehdotus, joka osoittaa kohteet, jotka haluat, muotoiltuna hakasulkeiden sisällä seuraavasti: "{[}Ehdotus: <numero> <kohteen nimi>, <numero> <kohteen nimi>, <...>{]}"} \\
\texttt{(d) Jos ehdotuksesi ovat täydentäviä, eli molempien ehdotusten täyttämiseksi on tarpeeksi kohteita, jokainen pelaaja saa pistemäärän saamiensa kohteiden arvojen summan perusteella. Muussa tapauksessa molemmat saavat nolla pistettä.} \\
\texttt{(e) Tavoitteesi on maksimoida pistemääräsi ja toisen pelaajan pistemäärän summa.} \\
\\ 
\texttt{Aloitetaan.} \\
\\ 
\texttt{Saatavilla olevien kohteiden joukko on:} \\
\\ 
\texttt{2 kynät, 2 sitruunat, 2 potit} \\
\\ 
\texttt{Salainen arvofunktiosi on:} \\
\\ 
\texttt{kynä: 3, sitruuna: 1, potti: 1} \\
\\ 
\texttt{TÄRKEÄÄ: Viestit, paitsi salaiset ehdotukset, lähetetään suoraan toiselle pelaajalle, joten älä sisällytä mitään vastausta sääntöihin tai tekstiä, joka ilmoittaa viestistäsi. Salaisesti ehdotuksen tekemiseksi käytä ilmoitettua muotoa. Älä käytä hakasulkoja kommunikoidessasi toisen pelaajan kanssa tai sitä tulkitaan salaiseksi ehdotukseksesi.} \\
\\ 
\texttt{Minulla on kynät ja potit arvokkaimpia. Ehdotan, että otan yhden potin ja yhden kynän, jos sinulle sopii. Miten sinä näet kohteiden arvot?} \\
            }
        }
    }
    & \\ \\

    \theutterance \stepcounter{utterance}  
    & & & \multicolumn{4}{p{0.6\linewidth}}{
        \cellcolor[rgb]{0.9,0.9,0.9}{
            \makecell[{{p{\linewidth}}}]{
                \texttt{\tiny{[GM$\langle$P2]}}
                \texttt{Minulle kynät ovat arvokkaimpia, sitruunat ja potit ovat samanarvoisia. Ottaisin mielelläni yhden kynän ja yhden sitruunan, jos se sopii sinulle.} \\
            }
        }
    }
     \\ \\

    \theutterance \stepcounter{utterance}  
    & & & \multicolumn{2}{p{0.3\linewidth}}{
        \cellcolor[rgb]{0.9,0.9,0.9}{
            \makecell[{{p{\linewidth}}}]{
                \texttt{\tiny{[GM$|$GM]}}
                \texttt{continue} \\
            }
        }
    }
    & & \\ \\

    \theutterance \stepcounter{utterance}  
    & & & \multicolumn{2}{p{0.3\linewidth}}{
        \cellcolor[rgb]{0.9,0.9,0.9}{
            \makecell[{{p{\linewidth}}}]{
                \texttt{\tiny{[GM$|$GM]}}
                \texttt{Minulle kynät ovat arvokkaimpia, sitruunat ja potit ovat samanarvoisia. Ottaisin mielelläni yhden kynän ja yhden sitruunan, jos se sopii sinulle.} \\
            }
        }
    }
    & & \\ \\

    \theutterance \stepcounter{utterance}  
    & & \multicolumn{4}{p{0.6\linewidth}}{
        \cellcolor[rgb]{0.9,0.9,0.9}{
            \makecell[{{p{\linewidth}}}]{
                \texttt{\tiny{[P1$\langle$GM]}}
                \texttt{Minulle kynät ovat arvokkaimpia, sitruunat ja potit ovat samanarvoisia. Ottaisin mielelläni yhden kynän ja yhden sitruunan, jos se sopii sinulle.} \\
            }
        }
    }
    & \\ \\

    \theutterance \stepcounter{utterance}  
    & \multicolumn{4}{p{0.6\linewidth}}{
        \cellcolor[rgb]{0.9,0.9,0.9}{
            \makecell[{{p{\linewidth}}}]{
                \texttt{\tiny{[P1$\rangle$GM]}}
                \texttt{Hyvä, voimmeko siis jakaa kohteet seuraavasti: minä otan yhden potin ja yhden kynän, ja sinä otat yhden kynän ja yhden sitruunan? Tämä jättää yhden potin ja yhden sitruunan jakamatta. Haluatko sinä ne vai minä?} \\
            }
        }
    }
    & & \\ \\

    \theutterance \stepcounter{utterance}  
    & & & \multicolumn{2}{p{0.3\linewidth}}{
        \cellcolor[rgb]{0.9,0.9,0.9}{
            \makecell[{{p{\linewidth}}}]{
                \texttt{\tiny{[GM$|$GM]}}
                \texttt{continue} \\
            }
        }
    }
    & & \\ \\

    \theutterance \stepcounter{utterance}  
    & & & \multicolumn{2}{p{0.3\linewidth}}{
        \cellcolor[rgb]{0.9,0.9,0.9}{
            \makecell[{{p{\linewidth}}}]{
                \texttt{\tiny{[GM$|$GM]}}
                \texttt{Hyvä, voimmeko siis jakaa kohteet seuraavasti: minä otan yhden potin ja yhden kynän, ja sinä otat yhden kynän ja yhden sitruunan? Tämä jättää yhden potin ja yhden sitruunan jakamatta. Haluatko sinä ne vai minä?} \\
            }
        }
    }
    & & \\ \\

    \theutterance \stepcounter{utterance}  
    & & \multicolumn{4}{p{0.6\linewidth}}{
        \cellcolor[rgb]{0.9,0.9,0.9}{
            \makecell[{{p{\linewidth}}}]{
                \texttt{\tiny{[GM$\rangle$P2]}}
                \texttt{Hyvä, voimmeko siis jakaa kohteet seuraavasti: minä otan yhden potin ja yhden kynän, ja sinä otat yhden kynän ja yhden sitruunan? Tämä jättää yhden potin ja yhden sitruunan jakamatta. Haluatko sinä ne vai minä?} \\
            }
        }
    }
    & \\ \\

    \theutterance \stepcounter{utterance}  
    & & & \multicolumn{4}{p{0.6\linewidth}}{
        \cellcolor[rgb]{0.9,0.9,0.9}{
            \makecell[{{p{\linewidth}}}]{
                \texttt{\tiny{[GM$\langle$P2]}}
                \texttt{Jos sinä otat jäljelle jäävän potin, minä otan sitruunan. Näin molemmat saavat kaksi kohdetta. Sopiiko tämä?} \\
            }
        }
    }
     \\ \\

    \theutterance \stepcounter{utterance}  
    & & & \multicolumn{2}{p{0.3\linewidth}}{
        \cellcolor[rgb]{0.9,0.9,0.9}{
            \makecell[{{p{\linewidth}}}]{
                \texttt{\tiny{[GM$|$GM]}}
                \texttt{continue} \\
            }
        }
    }
    & & \\ \\

    \theutterance \stepcounter{utterance}  
    & & & \multicolumn{2}{p{0.3\linewidth}}{
        \cellcolor[rgb]{0.9,0.9,0.9}{
            \makecell[{{p{\linewidth}}}]{
                \texttt{\tiny{[GM$|$GM]}}
                \texttt{Jos sinä otat jäljelle jäävän potin, minä otan sitruunan. Näin molemmat saavat kaksi kohdetta. Sopiiko tämä?} \\
            }
        }
    }
    & & \\ \\

    \theutterance \stepcounter{utterance}  
    & & \multicolumn{4}{p{0.6\linewidth}}{
        \cellcolor[rgb]{0.9,0.9,0.9}{
            \makecell[{{p{\linewidth}}}]{
                \texttt{\tiny{[P1$\langle$GM]}}
                \texttt{Jos sinä otat jäljelle jäävän potin, minä otan sitruunan. Näin molemmat saavat kaksi kohdetta. Sopiiko tämä?} \\
            }
        }
    }
    & \\ \\

    \theutterance \stepcounter{utterance}  
    & \multicolumn{4}{p{0.6\linewidth}}{
        \cellcolor[rgb]{0.9,0.9,0.9}{
            \makecell[{{p{\linewidth}}}]{
                \texttt{\tiny{[P1$\rangle$GM]}}
                \texttt{{[}Ehdotus: 1 kynä, 2 potit{]}} \\
            }
        }
    }
    & & \\ \\

    \theutterance \stepcounter{utterance}  
    & & & \multicolumn{2}{p{0.3\linewidth}}{
        \cellcolor[rgb]{0.9,0.9,0.9}{
            \makecell[{{p{\linewidth}}}]{
                \texttt{\tiny{[GM$|$GM]}}
                \texttt{proposal} \\
            }
        }
    }
    & & \\ \\

    \theutterance \stepcounter{utterance}  
    & & & \multicolumn{2}{p{0.3\linewidth}}{
        \cellcolor[rgb]{0.9,0.9,0.9}{
            \makecell[{{p{\linewidth}}}]{
                \texttt{\tiny{[GM$|$GM]}}
                [1, 0, 2]
            }
        }
    }
    & & \\ \\

    \theutterance \stepcounter{utterance}  
    & & \multicolumn{4}{p{0.6\linewidth}}{
        \cellcolor[rgb]{0.9,0.9,0.9}{
            \makecell[{{p{\linewidth}}}]{
                \texttt{\tiny{[GM$\rangle$P2]}}
                \texttt{Toinen pelaaja on jättänyt salaisen ehdotuksen. Nyt on sinun vuorosi jättää oma ehdotuksesi.} \\
            }
        }
    }
    & \\ \\

    \theutterance \stepcounter{utterance}  
    & & & \multicolumn{4}{p{0.6\linewidth}}{
        \cellcolor[rgb]{0.9,0.9,0.9}{
            \makecell[{{p{\linewidth}}}]{
                \texttt{\tiny{[GM$\langle$P2]}}
                \texttt{{[}Ehdotus: 1 kynä, 1 sitruuna{]}} \\
            }
        }
    }
     \\ \\

    \theutterance \stepcounter{utterance}  
    & & & \multicolumn{2}{p{0.3\linewidth}}{
        \cellcolor[rgb]{0.9,0.9,0.9}{
            \makecell[{{p{\linewidth}}}]{
                \texttt{\tiny{[GM$|$GM]}}
                \texttt{proposal} \\
            }
        }
    }
    & & \\ \\

    \theutterance \stepcounter{utterance}  
    & & & \multicolumn{2}{p{0.3\linewidth}}{
        \cellcolor[rgb]{0.9,0.9,0.9}{
            \makecell[{{p{\linewidth}}}]{
                \texttt{\tiny{[GM$|$GM]}}
                [1, 1, 0]
            }
        }
    }
    & & \\ \\

    \theutterance \stepcounter{utterance}  
    & & & \multicolumn{2}{p{0.3\linewidth}}{
        \cellcolor[rgb]{0.9,0.9,0.9}{
            \makecell[{{p{\linewidth}}}]{
                \texttt{\tiny{[GM$|$GM]}}
                [[1, 0, 2], [1, 1, 0]]
            }
        }
    }
    & & \\ \\

\end{supertabular}
}

\end{document}
